% called by main.tex
%
\chapter{Material y métodos}
\label{sec::capitulo3}

Este capítulo describe los materiales y la metodología empleados en el trabajo:
el conjunto de datos, la cadena de procesamiento (ETL), las configuraciones de
clases y el diseño experimental, la arquitectura de los modelos, el protocolo de
entrenamiento y evaluación, la técnica de explicabilidad y los mecanismos de
gestión de experimentos y reproducibilidad. La descripción se corresponde con la
implementación desarrollada en el repositorio del proyecto.


\section{Conjunto de datos: CheXpert}
\label{sec::mm_dataset}

El trabajo se basa en \textbf{CheXpert}~\cite{irvinCheXpertLargeChest2019}, un
conjunto de datos público de la Universidad de Stanford que contiene 224.316
radiografías de tórax de 65.240 pacientes, recogidas de forma retrospectiva. Cada
estudio está etiquetado con 14 observaciones clínicas extraídas automáticamente
de los informes radiológicos mediante procesamiento del lenguaje natural. Para
cada observación, la etiqueta puede tomar el valor positivo (1), negativo (0),
ausente (sin mención) o incierto ($-1$), siendo este último uno de los rasgos
distintivos del conjunto.

En este trabajo se emplean \textbf{13 de las 14 observaciones}, excluyendo
\textit{Pleural Other} por su muy baja prevalencia en el conjunto de
entrenamiento, que no aporta señal estadística suficiente. La clase
\textit{No Finding} se conserva para que el modelo pueda indicar explícitamente
la ausencia de hallazgos. Como conjunto de prueba se utiliza el conjunto de
validación oficial de CheXpert, anotado por consenso de radiólogos de Stanford;
al estar verificado por expertos, actúa como referencia \textit{silver-standard}
y es independiente del entrenamiento (pacientes disjuntos), por lo que no
participa en la selección del modelo.


\section{Procesamiento de datos (ETL)}
\label{sec::mm_etl}

El conjunto de datos se somete a una cadena de extracción, transformación y carga
(ETL) implementada como una lista blanca (\textit{whitelist}) de criterios, lo
que estandariza el preprocesado y evita sesgos de selección arbitrarios. Los
criterios de retención son:

\begin{itemize}
  \item \textbf{Proyección.} Solo se conservan las radiografías de vista frontal
        y proyección anteroposterior (AP), por coherencia con el diseño del
        trabajo.
  \item \textbf{Valores de etiqueta permitidos.} Se admiten únicamente los
        valores $\{0,\ 1,\ \text{ausente}\}$. Las etiquetas inciertas ($-1$) se
        descartan, lo que corresponde a la estrategia de gestión de la
        incertidumbre \textbf{U-Ignore} (véase la Sección~\ref{sec::mm_clases}).
  \item \textbf{Coherencia de \textit{No Finding}.} Se eliminan los estudios
        inconsistentes en los que \textit{No Finding} aparece como positivo
        simultáneamente con alguna otra patología.
  \item \textbf{Sexo.} Se conservan los registros de pacientes de sexo masculino
        y femenino.
\end{itemize}

Tras el filtrado, las etiquetas ausentes se imputan a 0 (negativo), ya que la
ausencia de mención en el informe constituye la señal más débil posible,
asimilable a un negativo implícito. La división en entrenamiento y validación se
realiza \textbf{a nivel de paciente} (90\,\% / 10\,\%), garantizando que todas
las imágenes de un mismo paciente pertenezcan a un único subconjunto. Esta
partición por paciente es esencial para evitar la fuga de información
(\textit{data leakage}): si imágenes del mismo paciente apareciesen en
entrenamiento y validación, el modelo podría aprender rasgos individuales de
estilo de adquisición en lugar de patrones clínicos generalizables, produciendo
una estimación optimista del rendimiento.


\section{Configuraciones de clases y diseño experimental}
\label{sec::mm_clases}

\subsection{Gestión de la incertidumbre}

CheXpert codifica explícitamente la incertidumbre de ciertas etiquetas con el
valor $-1$. En este trabajo se adopta la estrategia \textbf{U-Ignore}, que
consiste en descartar las muestras con etiquetas inciertas, implementada en el
ETL mediante la lista blanca de valores permitidos. Esta estrategia es la más
conservadora —no introduce supuestos sobre el significado de la incertidumbre— y
sirve como base metodológica del trabajo. La evaluación comparativa de las demás
estrategias descritas en la literatura (U-Zeros, U-Ones, U-SelfTrained y
U-MultiClass)~\cite{irvinCheXpertLargeChest2019} queda planteada como línea de
trabajo futuro.

\subsection{Configuraciones de clases}

El sistema permite seleccionar, mediante configuración, qué conjunto de clases
activa el modelo. Cada configuración fija sus clases activas y una política de
eliminación de estudios (\textit{anti-ruido}) que descarta los estudios cuyas
únicas etiquetas positivas pertenecen a clases eliminadas, evitando convertir un
positivo real en un falso negativo. Se definen tres configuraciones:

\begin{table}[h]
\centering
\caption{Configuraciones de clases evaluadas en el trabajo.}
\label{tab::mm_configs}
\begin{tabularx}{\textwidth}{l c X}
\toprule
\textbf{Configuración} & \textbf{N.º clases} & \textbf{Descripción} \\
\midrule
\texttt{full13} & 13 & Todas las clases activas (sin \textit{Pleural Other}). Sin eliminación de estudios. \\
\texttt{nofracture12} & 12 & Elimina \textit{Fracture}. Descarta los estudios huérfanos cuya única etiqueta positiva era \textit{Fracture}. \\
\texttt{min5pct9} & 9 & Elimina las clases de prevalencia inferior al 5\,\% (\textit{Enlarged Cardiomediastinum}, \textit{Lung Lesion}, \textit{Pneumonia} y \textit{Fracture}). Descarta los estudios sin ninguna etiqueta positiva en las clases activas. \\
\bottomrule
\end{tabularx}
\end{table}

\subsection{Diseño experimental}

El estudio se organiza en torno a un conjunto de experimentos que combinan las
configuraciones de clases anteriores con distintas arquitecturas. La comparación
de configuraciones de clases se realiza con una arquitectura de referencia
(DenseNet-121), mientras que la comparación de arquitecturas se realiza sobre una
configuración de clases fija; no se evalúa el producto cartesiano completo de
arquitecturas y configuraciones por restricciones de tiempo de cómputo. Los
experimentos planificados son los siguientes:

\textbf{Experimento 1 — DenseNet-121 con 13 clases (\texttt{full13}).}
\begin{itemize}
  \item \textit{Objetivo:} establecer un modelo de referencia sobre el conjunto
        completo de 13 clases.
  \item \textit{Configuración:} \texttt{full13} (sin \textit{Pleural Other}).
  \item \textit{Motivo:} constituye la referencia más amplia y comparable con la
        literatura. Este experimento debe repetirse para verificar la
        consistencia de los resultados.
  \item \textit{Ventajas:} máxima cobertura diagnóstica. \textit{Inconvenientes:}
        incluye clases muy poco frecuentes y difíciles de evaluar de forma
        fiable.
\end{itemize}

\textbf{Experimento 2 — DenseNet-121 con 12 clases (\texttt{nofracture12}).}
\begin{itemize}
  \item \textit{Objetivo:} entrenar y evaluar sin la clase \textit{Fracture}.
  \item \textit{Configuración:} \texttt{nofracture12}.
  \item \textit{Motivo:} se detectó que \textit{Fracture} no dispone de muestras
        positivas en el conjunto de prueba, por lo que no puede evaluarse de
        forma correcta; su eliminación produce una evaluación más fiable.
  \item \textit{Ventajas:} métricas de prueba más representativas.
        \textit{Inconvenientes:} se pierde la capacidad de detectar fracturas.
\end{itemize}

\textbf{Experimento 3 — DenseNet-121 con 9 clases (\texttt{min5pct9}, pendiente).}
\begin{itemize}
  \item \textit{Objetivo:} analizar el comportamiento en un escenario menos
        desbalanceado, conservando únicamente las clases con prevalencia igual o
        superior al 5\,\%.
  \item \textit{Configuración:} \texttt{min5pct9} (9 clases).
  \item \textit{Motivo:} reducir el desbalanceo extremo eliminando las clases
        menos representadas, lo que puede dar lugar a un modelo más estable.
  \item \textit{Ventajas:} menor desbalanceo y mayor representatividad de las
        clases. \textit{Inconvenientes:} menor cobertura diagnóstica; queda
        supeditado a la disponibilidad de tiempo de cómputo.
\end{itemize}

Sobre una configuración de referencia se evalúan de forma comparativa varias
arquitecturas (Sección~\ref{sec::mm_arquitectura}); el alcance final de esta
comparativa dependerá del tiempo de cómputo disponible.


\section{Arquitectura del modelo}
\label{sec::mm_arquitectura}

El modelo de clasificación sigue un esquema de \textbf{aprendizaje por
transferencia}: un \textit{backbone} convolucional preentrenado sobre
ImageNet~\cite{dengImageNet2009} actúa como extractor de características, y su
capa de clasificación original se reemplaza por una cabeza adaptada a la tarea
multietiqueta. La cabeza de clasificación, común a todos los backbones, se
compone de cuatro capas:

\begin{equation*}
\text{Linear}(d \rightarrow 1024) \;\rightarrow\; \text{ReLU} \;\rightarrow\;
\text{Dropout}(0{,}5) \;\rightarrow\; \text{Linear}(1024 \rightarrow C),
\end{equation*}

donde $d$ es la dimensión de características del backbone y $C$ el número de
clases activas. La capa oculta permite aprender representaciones no lineales de
alto nivel, mientras que la regularización por
\textit{Dropout}~\cite{srivastavaDropout2014} mitiga el sobreajuste, especialmente
relevante en imagen médica con muestras limitadas. La salida son $C$ logits, uno
por patología, sobre los que se aplica una función sigmoide para obtener
probabilidades independientes.

Las radiografías, originalmente en escala de grises, se convierten a tres canales
(replicando el canal único) para ser compatibles con los pesos preentrenados en
ImageNet. Se evalúan cinco arquitecturas de referencia, intercambiables mediante
configuración:
\textbf{DenseNet-121}~\cite{huangDenselyConnectedConvolutional2017},
\textbf{ResNet-50}~\cite{heDeepResidualLearning2016},
\textbf{VGG16-BN}~\cite{simonyanVeryDeepConvolutional2015},
\textbf{ConvNeXt-Tiny}~\cite{liuConvNet2020s2022} y
\textbf{Swin-T}~\cite{liuSwinTransformer2021}, abarcando desde CNN clásicas hasta
transformers de visión.


\section{Función de pérdida y protocolo de entrenamiento}
\label{sec::mm_entrenamiento}

\subsection{Función de pérdida}

Al tratarse de un problema multietiqueta, se emplea la entropía cruzada binaria
sobre cada salida independiente. Para contrarrestar el fuerte desbalanceo de
clases, se utiliza \texttt{BCEWithLogitsLoss} con un peso por clase
$\text{pos\_weight} = n_{\text{neg}} / n_{\text{pos}}$, calculado sobre el
conjunto de entrenamiento. Este peso escala el gradiente de los ejemplos
positivos de cada patología en proporción inversa a su frecuencia, de modo que
las clases minoritarias reciban una atención comparable a las
mayoritarias~\cite{usharubyBinaryCrossEntropy2020}.

\subsection{Optimización y regularización}

Los hiperparámetros de entrenamiento se resumen en la
Tabla~\ref{tab::mm_hiperparametros}. El modelo se optimiza con
\textbf{AdamW}~\cite{kingmaAdamMethod2015}, que incorpora regularización por
decaimiento de pesos (\textit{weight decay}). La tasa de aprendizaje se ajusta
dinámicamente con el planificador \texttt{ReduceLROnPlateau}, que la reduce
cuando la pérdida de validación deja de mejorar. El entrenamiento emplea
precisión mixta automática (AMP) en GPU para reducir el consumo de memoria y
acelerar el cómputo. Se aplica \textbf{parada temprana} (\textit{early stopping})
cuando la pérdida de validación no mejora durante un número fijo de épocas,
restaurándose al final los pesos del epoch con mejor AUROC de validación.

\begin{table}[h]
\centering
\caption{Hiperparámetros de entrenamiento.}
\label{tab::mm_hiperparametros}
\begin{tabularx}{\textwidth}{l l X}
\toprule
\textbf{Hiperparámetro} & \textbf{Valor} & \textbf{Descripción} \\
\midrule
Optimizador & AdamW & Tasa de aprendizaje $10^{-4}$, \textit{weight decay} $0{,}01$ \\
Planificador & ReduceLROnPlateau & Paciencia 2 épocas \\
Tamaño de lote & 64 / 256 & Entrenamiento / inferencia \\
Épocas máximas & 50 & Con parada temprana (paciencia 8) \\
Tamaño de imagen & $224 \times 224$ & Píxeles \\
Precisión mixta & AMP & Activada en GPU \\
Semilla & 42 & Reproducibilidad \\
\bottomrule
\end{tabularx}
\end{table}

\subsection{Aumento de datos}

Sobre el conjunto de entrenamiento se aplica aumento de datos para incrementar la
variabilidad y reducir el sobreajuste: volteo horizontal aleatorio (probabilidad
0,5) y una transformación afín aleatoria (rotación de hasta $10°$ y traslación de
hasta el 5\,\%). Estas transformaciones simulan la variabilidad natural de
posición del paciente sin introducir artefactos clínicamente incorrectos. Tras el
aumento, las imágenes se redimensionan a $224 \times 224$ y se normalizan con la
media y desviación estándar de ImageNet. En validación y prueba no se aplica
aumento, para obtener una estimación determinista del rendimiento.


\section{Evaluación y métricas}
\label{sec::mm_metricas}

La evaluación se realiza en dos niveles. En \textbf{validación}, durante el
entrenamiento, se monitoriza el AUROC-macro (área bajo la curva ROC promediada
sobre las clases evaluables), que es independiente del umbral de decisión y
robusto frente al desbalanceo; este criterio determina la selección del mejor
epoch. En el conjunto de \textbf{prueba} \textit{silver-standard} se calcula un
conjunto más amplio de métricas por patología y agregadas: AUROC, PR-AUC (área
bajo la curva de precisión-exhaustividad), precisión, exhaustividad y F1, además
del AUROC promediado sobre las cinco patologías de la competición oficial de
CheXpert (\textit{Atelectasis}, \textit{Cardiomegaly}, \textit{Consolidation},
\textit{Edema} y \textit{Pleural Effusion}), bien representadas en el conjunto de
prueba.

La selección del modelo de producción se realiza mediante un \textbf{mecanismo de
promoción} automático: un modelo recién entrenado solo reemplaza al mejor modelo
registrado para su mismo par (arquitectura, configuración de clases) si mejora el
AUROC de las cinco patologías oficiales sobre el conjunto de prueba por encima de
un margen mínimo (0,005), empleándose el F1-macro como criterio de desempate.
Este margen evita sustituciones por ruido estadístico. Las predicciones binarias,
cuando se requieren, se obtienen con un umbral de 0,5.


\section{Explicabilidad mediante Grad-CAM}
\label{sec::mm_gradcam}

Para dotar al sistema de capacidad explicativa se emplea
\textbf{Grad-CAM}~\cite{selvarajuGradCAMVisualExplanations2017}, que genera mapas
de calor a partir de los gradientes que fluyen hacia la última capa convolucional
del modelo, señalando las regiones de la imagen más influyentes en la predicción
de cada patología. La capa objetivo se selecciona de forma específica para cada
arquitectura (la última etapa convolucional del backbone); en el caso del Swin
Transformer, cuya salida conserva los tokens en formato \textit{channels-last},
se aplica una reorganización de ejes para adaptarla al formato espacial que
espera Grad-CAM. Esta técnica no requiere modificar la arquitectura ni reentrenar
el modelo. Las explicaciones generadas se integran en la plataforma web
(Capítulo~\ref{sec::capitulo5}) y se interpretan como una ayuda a la verificación
por parte del especialista, no como una garantía de corrección del modelo, en
línea con las cautelas señaladas en la literatura sobre la fiabilidad de los
mapas de saliencia~\cite{zhangRevisitingTrustworthiness2024}.


\section{Gestión de experimentos y reproducibilidad}
\label{sec::mm_reproducibilidad}

Cada ejecución de entrenamiento se documenta de forma automática y trazable. El
sistema registra, por cada experimento, los metadatos de ejecución (identificador
del \textit{commit} de git, hardware, duración, mejor epoch), la configuración
empleada, la distribución de clases de los conjuntos, el historial de métricas
por época, las métricas finales de validación y prueba por patología, las
predicciones almacenadas, un análisis de los peores casos (falsos positivos y
negativos) y las figuras de evaluación. Un registro centralizado mantiene el
mejor modelo por cada par (arquitectura, configuración de clases) y un índice
acumulado (\textit{leaderboard}) permite comparar las ejecuciones.

La reproducibilidad se garantiza fijando una semilla global única que controla
todas las fuentes de aleatoriedad (la variable \texttt{PYTHONHASHSEED}, los
generadores de NumPy, del módulo \texttt{random} de Python y de PyTorch en CPU y
GPU) y activando el modo determinista de cuDNN. El entorno de ejecución, basado
en PyTorch~\cite{paszkePyTorch2019}, y el código completo se documentan en el
repositorio del proyecto (Capítulo~\ref{sec::capitulo7}).
